\documentclass[10pt,a4paper]{amsart}
\usepackage{amsmath,amssymb,amsthm,graphicx}
% Retained preamble for the Stacks Project chapter snapshot in this directory.
% The chapter's own preamble.tex is a copy of the shared Stacks Project preamble; the block
% below is sliced MECHANICALLY, line for line, from it. Nothing was added, removed or
% redefined; the wrapper class, the xr-hyper/hyperref wiring and the externaldocument list
% are omitted because the assembler supplies the document class.
%
% lines taken from preamble.tex: 3, 5, 9, 11, 12, 156, 157, 158, ... (51 lines total)
%
% package coverage restored from the chapter's own preamble (2026-10-13): xy, xy-2cell

\usepackage[all]{xy}
\xyoption{2cell}
\UseAllTwocells

\newenvironment{reference}{\comment}{\endcomment}
\newenvironment{slogan}{\comment}{\endcomment}
\newenvironment{history}{\comment}{\endcomment}
\theoremstyle{plain}
\newtheorem{theorem}[subsection]{Theorem}
\newtheorem{proposition}[subsection]{Proposition}
\newtheorem{lemma}[subsection]{Lemma}
\theoremstyle{definition}
\newtheorem{definition}[subsection]{Definition}
\newtheorem{example}[subsection]{Example}
\newtheorem{exercise}[subsection]{Exercise}
\newtheorem{situation}[subsection]{Situation}
\theoremstyle{remark}
\newtheorem{remark}[subsection]{Remark}
\newtheorem{remarks}[subsection]{Remarks}
\numberwithin{equation}{subsection}
\def\lim{\mathop{\mathrm{lim}}\nolimits}
\def\colim{\mathop{\mathrm{colim}}\nolimits}
\def\Spec{\mathop{\mathrm{Spec}}}
\def\Hom{\mathop{\mathrm{Hom}}\nolimits}
\def\Ext{\mathop{\mathrm{Ext}}\nolimits}
\def\SheafHom{\mathop{\mathcal{H}\!\mathit{om}}\nolimits}
\def\SheafExt{\mathop{\mathcal{E}\!\mathit{xt}}\nolimits}
\def\Sch{\mathit{Sch}}
\def\Mor{\mathop{\mathrm{Mor}}\nolimits}
\def\Ob{\mathop{\mathrm{Ob}}\nolimits}
\def\Sh{\mathop{\mathit{Sh}}\nolimits}
\def\NL{\mathop{N\!L}\nolimits}
\def\CH{\mathop{\mathrm{CH}}\nolimits}
\def\proetale{{pro\text{-}\acute{e}tale}}
\def\etale{{\acute{e}tale}}
\def\QCoh{\mathit{QCoh}}
\def\Ker{\mathop{\mathrm{Ker}}}
\def\Im{\mathop{\mathrm{Im}}}
\def\Coker{\mathop{\mathrm{Coker}}}
\def\Coim{\mathop{\mathrm{Coim}}}
\DeclareMathSymbol{\boxtimes}{\mathbin}{AMSa}{"02}
\def\QCohstack{\mathcal{QC}\!\mathit{oh}}
\def\Cohstack{\mathcal{C}\!\mathit{oh}}
\def\Spacesstack{\mathcal{S}\!\mathit{paces}}
\def\Quotfunctor{\mathrm{Quot}}
\def\Hilbfunctor{\mathrm{Hilb}}
\def\Curvesstack{\mathcal{C}\!\mathit{urves}}
\def\Polarizedstack{\mathcal{P}\!\mathit{olarized}}
\def\Complexesstack{\mathcal{C}\!\mathit{omplexes}}
\def\Pic{\mathop{\mathrm{Pic}}\nolimits}
\def\Picardstack{\mathcal{P}\!\mathit{ic}}
\def\Picardfunctor{\mathrm{Pic}}
\def\Deformationcategory{\mathcal{D}\!\mathit{ef}}
% lane-s11: discard bibliographic \begin{reference} blocks entirely (incl. any \cite inside).
\newsavebox{\limrefbox}
\ifcsname reference\endcsname\expandafter\let\csname reference\endcsname\undefined\fi\ifcsname endreference\endcsname\expandafter\let\csname endreference\endcsname\undefined\fi
\def\citation#1{\relax}
\newenvironment{reference}{\setbox\limrefbox=\vbox\bgroup\let\cite\relax}{\egroup\ignorespaces}
% lane-s11: the chapter's retained preamble defines slogan/history via comment.sty's
% \comment macro, which is never loaded; redefine both with the same guarded
% discard pattern (content is authorial commentary, kept verbatim in the spans).
\ifcsname slogan\endcsname\expandafter\let\csname slogan\endcsname\undefined\fi\ifcsname endslogan\endcsname\expandafter\let\csname endslogan\endcsname\undefined\fi
\newenvironment{slogan}{\setbox\limrefbox=\vbox\bgroup}{\egroup\ignorespaces}
\ifcsname history\endcsname\expandafter\let\csname history\endcsname\undefined\fi\ifcsname endhistory\endcsname\expandafter\let\csname endhistory\endcsname\undefined\fi
\newenvironment{history}{\setbox\limrefbox=\vbox\bgroup}{\egroup\ignorespaces}
\usepackage{cleveref}
\crefname{originalreference}{original reference}{original references}
\Crefname{originalreference}{Original reference}{Original references}
\makeatletter
\newlabel{situation-descent}{{086P}{source}{}{sourcepointer-1}{}}
\newlabel{situation-descent@cref}{{[originalreference][0][]086P}{[1][0][]source}{}{}{}}
\makeatother
\crefname{originallemma}{original lemma}{original lemmas}
\Crefname{originallemma}{Original lemma}{Original lemmas}
\makeatletter
\newlabel{lemma-limit-affine}{{01Z6}{source}{}{sourcepointer-2}{}}
\newlabel{lemma-limit-affine@cref}{{[originallemma][0][]01Z6}{[1][0][]source}{}{}{}}
\makeatother
\crefname{originalreference}{original reference}{original references}
\Crefname{originalreference}{Original reference}{Original references}
\makeatletter
\newlabel{schemes-lemma-characterize-separated}{{01KP}{source}{}{sourcepointer-3}{}}
\newlabel{schemes-lemma-characterize-separated@cref}{{[originalreference][0][]01KP}{[1][0][]source}{}{}{}}
\makeatother
\crefname{originalreference}{original reference}{original references}
\Crefname{originalreference}{Original reference}{Original references}
\makeatletter
\newlabel{schemes-lemma-diagonal-immersion}{{01KJ}{source}{}{sourcepointer-4}{}}
\newlabel{schemes-lemma-diagonal-immersion@cref}{{[originalreference][0][]01KJ}{[1][0][]source}{}{}{}}
\makeatother
\crefname{originalreference}{original reference}{original references}
\Crefname{originalreference}{Original reference}{Original references}
\makeatletter
\newlabel{morphisms-lemma-locally-finite-type-characterize}{{01T2}{source}{}{sourcepointer-5}{}}
\newlabel{morphisms-lemma-locally-finite-type-characterize@cref}{{[originalreference][0][]01T2}{[1][0][]source}{}{}{}}
\makeatother
\crefname{originalreference}{original reference}{original references}
\Crefname{originalreference}{Original reference}{Original references}
\makeatletter
\newlabel{morphisms-lemma-immersion-locally-finite-type}{{01T5}{source}{}{sourcepointer-6}{}}
\newlabel{morphisms-lemma-immersion-locally-finite-type@cref}{{[originalreference][0][]01T5}{[1][0][]source}{}{}{}}
\makeatother
\title{Separatedness of a limit of schemes with separated transition maps}
\author{The Stacks Project Authors}
\date{stacks-project:limits.tex:local-snapshot-20261012; GNU FDL 1.2 or later}
\begin{document}
\maketitle
\noindent\textit{Editorial note.} Source typo kept verbatim: 'O(U \$\textbackslash{}otimes\$ V)' where the argument uses O(U \$\textbackslash{}cap\$ V) (tensor for intersection); recorded, not corrected in the counted span.\par\medskip
\noindent\textit{Allowed standard prerequisites (editorial).}\begin{itemize}
\item Generators of the limit topology. Limits, Section on properties of limits
\end{itemize}\medskip
\ifdefined\hypertarget\hypertarget{sourcepointer-1}{}\fi
\noindent\textit{Source reference.} Original reference 086P: sources/stacks-chow/tags.txt, source lines 10399--10399. Cross-chapter standard prerequisite cited by the author inside the counted span; it is used as a pointer and its own proof is not part of this item.\par\medskip
\ifdefined\hypertarget\hypertarget{sourcepointer-2}{}\fi
\noindent\textit{Source reference.} Original lemma 01Z6: sources/stacks-chow/tags.txt, source lines 2581--2581. Cross-chapter standard prerequisite cited by the author inside the counted span; it is used as a pointer and its own proof is not part of this item.\par\medskip
\ifdefined\hypertarget\hypertarget{sourcepointer-3}{}\fi
\noindent\textit{Source reference.} Original reference 01KP: sources/stacks-chow/tags.txt, source lines 2058--2058. Cross-chapter standard prerequisite cited by the author inside the counted span; it is used as a pointer and its own proof is not part of this item.\par\medskip
\ifdefined\hypertarget\hypertarget{sourcepointer-4}{}\fi
\noindent\textit{Source reference.} Original reference 01KJ: sources/stacks-chow/tags.txt, source lines 2052--2052. Cross-chapter standard prerequisite cited by the author inside the counted span; it is used as a pointer and its own proof is not part of this item.\par\medskip
\ifdefined\hypertarget\hypertarget{sourcepointer-5}{}\fi
\noindent\textit{Source reference.} Original reference 01T2: sources/stacks-chow/tags.txt, source lines 2359--2359. Cross-chapter standard prerequisite cited by the author inside the counted span; it is used as a pointer and its own proof is not part of this item.\par\medskip
\ifdefined\hypertarget\hypertarget{sourcepointer-6}{}\fi
\noindent\textit{Source reference.} Original reference 01T5: sources/stacks-chow/tags.txt, source lines 2362--2362. Cross-chapter standard prerequisite cited by the author inside the counted span; it is used as a pointer and its own proof is not part of this item.\par\medskip
\begin{lemma}
\label{lemma-limit-separated}
In Situation \ref{situation-descent} if $S$ is separated,
then for some $i_0 \in I$ the schemes $S_i$ for $i \geq i_0$
are separated.
\end{lemma}

\begin{proof}
Choose a finite affine open covering
$S_0 = U_{0, 1} \cup \ldots \cup U_{0, m}$.
Set $U_{i, j} \subset S_i$ and $U_j \subset S$
equal to the inverse image of $U_{0, j}$.
Note that $U_{i, j}$ and $U_j$ are affine. As $S$ is separated
the intersections $U_{j_1} \cap U_{j_2}$ are affine. Since
$U_{j_1} \cap U_{j_2} = \lim_{i \geq 0} U_{i, j_1} \cap U_{i, j_2}$
we see that $U_{i, j_1} \cap U_{i, j_2}$ is affine for large $i$
by Lemma \ref{lemma-limit-affine}. To show that $S_i$ is separated
for large $i$ it now suffices to show that
$$
\mathcal{O}_{S_i}(U_{i, j_1})
\otimes_{\mathcal{O}_S(S)}
\mathcal{O}_{S_i}(U_{i, j_2})
\longrightarrow
\mathcal{O}_{S_i}(U_{i, j_1} \cap U_{i, j_2})
$$
is surjective for large $i$
(Schemes, Lemma \ref{schemes-lemma-characterize-separated}).

\medskip\noindent
To get rid of the annoying indices, assume we have affine opens
$U, V \subset S_0$ such that $U \cap V$ is affine too.
Let $U_i, V_i \subset S_i$, resp.\ $U, V \subset S$ be the inverse images.
We have to show that
$\mathcal{O}(U_i) \otimes \mathcal{O}(V_i) \to
\mathcal{O}(U_i \cap V_i)$
is surjective for $i$ large enough and we know that
$\mathcal{O}(U) \otimes \mathcal{O}(V) \to \mathcal{O}(U \cap V)$
is surjective. Note that
$\mathcal{O}(U_0) \otimes \mathcal{O}(V_0) \to
\mathcal{O}(U_0 \cap V_0)$
is of finite type, as the diagonal morphism $S_i \to S_i \times S_i$
is an immersion (Schemes, Lemma \ref{schemes-lemma-diagonal-immersion})
hence locally of finite type
(Morphisms, Lemmas \ref{morphisms-lemma-locally-finite-type-characterize} and
\ref{morphisms-lemma-immersion-locally-finite-type}).
Thus we can choose elements
$f_{0, 1}, \ldots, f_{0, n} \in \mathcal{O}(U_0 \cap V_0)$
which generate $\mathcal{O}(U_0 \cap V_0)$ over
$\mathcal{O}(U_0) \otimes \mathcal{O}(V_0)$.
Observe that for $i \geq 0$ the diagram of schemes
$$
\xymatrix{
U_i \cap V_i \ar[r] \ar[d] & U_i \ar[d] \\
U_0 \cap V_0 \ar[r] & U_0
}
$$
is cartesian. Thus we see that the images
$f_{i, 1}, \ldots, f_{i, n} \in \mathcal{O}(U_i \cap V_i)$
generate $\mathcal{O}(U_i \cap V_i)$ over
$\mathcal{O}(U_i) \otimes \mathcal{O}(V_0)$
and a fortiori over
$\mathcal{O}(U_i) \otimes \mathcal{O}(V_i)$.
By assumption the images $f_1, \ldots, f_n \in \mathcal{O}(U \otimes V)$
are in the image of the map
$\mathcal{O}(U) \otimes \mathcal{O}(V) \to \mathcal{O}(U \cap V)$.
Since
$\mathcal{O}(U) \otimes \mathcal{O}(V) =
\colim \mathcal{O}(U_i) \otimes \mathcal{O}(V_i)$
we see that they are in the image of the map at some finite level
and the lemma is proved.
\end{proof}

\end{document}
